% Part: lambda-calculus
% Chapter: lambda-definability
% Section: partial-recursive-functions

\documentclass[../../../include/open-logic-section]{subfiles}

\begin{document}

\olfileid{lam}{ldf}{par}
\olsection{جزوي بازګشتي تابعې \usetoken{S}{lambda definable} دي}

جزوي بازګشتي تابعې له بنسټيزو تابعو څخه د ترکيب،
بنسټيز بازګښت او نامحدودې کمينه‌موندنې په وسيله
ترلاسه کېږي۔ له عامو بازګشتي تابعو سره ئې توپير
دا دے چې په بې‌حده پلټنه کښې کارېدونکې تابعې
لازم نۀ دي چې منظمې وي۔ د منظم‌والي د شرط نشتوالے
په دې معنا دے چې د کمينه‌موندنې په وسيله تعريف شوې
تابعه کله ناکله ناتعريفه کېدے شي۔

په لومړي نظر ښايي داسې وبرېښي چې د ټولو هرځاے
تعريف شوو عامو بازګشتي تابعو د !!{lambda definable}
والي د ثبوت هماغه لارې د ټولو جزوي بازګشتي تابعو
د !!{lambda definable} والي د ثبوت دپاره هم کارېدے شي۔
د بېلګې په توګه، که $f$ او $g$ په ترتيب سره د $F$
او~$G$ ترمونو په وسيله !!{lambda define}d وي،
نو د $f$ او~$g$ ترکيب د
$\lambd[x][F (G x)]$ په وسيله !!{lambda define}d
دے۔ خو کله چې تابعې جزوي وي، ستونزه راپيدا کېږي۔
که $g(x)$ ناتعريفه وي، د $G x$ عادي بڼه نۀ وي۔
په ډېرو حالاتو کښې له دې امله $F (G x)$ هم عادي
بڼه نۀ لري؛ همدا مو غوښتل۔ خو که $F$ د
$\lambd[x][\lambd[y][y]]$ ترم وي، بيا $F (G x)$
عادي بڼه لري: $\lambd[y][y]$۔

دا ستونزه د حل وړ ده۔ د ټولو جزوي بازګشتي تابعو
د !!{lambda define} داسې لارې شته چې ناتعريفه
ارزښتونه ئې په هغو ترمونو تمثيلېږي چې عادي بڼه نۀ لري۔ خو دا
لارې د هغو لارو په پرتله چې د عامو بازګشتي تابعو
دپاره مو کارولې، څه پېچلې او لږې شهودي دي۔ دا قضيه
دلته بې له ثبوته ثبتوو:

\begin{thm}
  ټولې جزوي بازګشتي تابعې !!{lambda definable} دي۔
\end{thm}

\end{document}
